\chapter{DRIVE 视网膜血管分割}
\label{chap:drive}

\section{数据与任务}
输入为彩色眼底照片，输出为血管像素组成的二值掩膜。课程目录含
20 张带血管标注的训练图和 20 张用于提交的测试图，图像来自 DRIVE
基准\citep{staal2004drive}。血管占全图像素比例较低，细血管和分支
与背景对比度低。只用像素准确率会掩盖对血管前景的漏检，因此课程
以 Dice 评价。官方测试标注以及由它衍生的掩膜不进入主线的训练、
阈值选择或效果展示。选定模型使用可追溯来源的公开权重和课程
图像推理。

设 $P,G$ 分别为预测和标注的血管像素集合，
\begin{equation}
  D(P,G)=\frac{2|P\cap G|}{|P|+|G|}
  =\frac{2TP}{2TP+FP+FN},\qquad
  \operatorname{IoU}=\frac{TP}{TP+FP+FN}.
  \label{eq:drive-dice}
\end{equation}
分割代码可逐图计算再求均值，原生 U-Net 训练器将验证图的像素计数累加后
计算 F1。二者的权重口径不同。是否将视野外预测清零也会改变假阳性计数。
原 DRIVE 论文的 AUC 和准确率适用于原基准评价，本文的课程榜单使用实际
提交得到的 Dice。

\section{模型与推理}
\subsection{FR-UNet 多分辨率特征交换}
FR-UNet 在整个网络中维持全分辨率分支，同时引入较低分辨率分支扩大
感受范围\citep{liu2022frunet}。本地适配输入为 $B\times1\times512\times512$。
尺度系数为 2，对应实现中的四级主分支宽度 32、64、128 和 256。
各级通过下采样、上采样和横向残差块交换特征，使用拼接保留来源信息。
最终从五个全分辨率阶段分别经 $1\times1$ 卷积产生 logits，
按层平均后再做 Sigmoid。本地代码返回 logits，使推理概率变换统一在
模型调用层完成。高分辨率分支有利于保留细线结构，密集跨尺度连接也增加了
中间特征存储。

\archdiagram{drive_arch}{冻结 FR-UNet 的推理与提交路径}{fig:drive-arch}

所选 checkpoint 是上游 DRIVE 第 40 轮权重的本地封装，许可为 MIT。
第 40 轮来自上游模型标识。课程最终路径直接加载该文件，使用单尺度
512 缩放和单次推理。原始训练的硬件和完整逐轮日志在本地材料中不可查，
实际自训的训练过程另列于第~\ref{sec:native-train} 节。

\subsection{双阈值连通与提交编码}
高阈值产生种子 $M_0=\mathbf 1[p\ge t_h]$，低阈值定义可扩张区域
$A=\mathbf 1[p\ge t_l]$。令 $\operatorname{Dilate}_{3\times3}$ 表示
采用 $3\times3$ 邻域的膨胀，迭代计算
\[
M_{k+1}=\operatorname{Dilate}_{3\times3}(M_k)\cap A .
\]
当 $M_{k+1}=M_k$ 时停止。由于 $t_l\le t_h$ 且膨胀包含中心像素，
该序列单调扩张并在有限网格内收敛。选定配方使用 $t_h=0.65$、
$t_l=0.40$。处理顺序为双阈值迭代、与估计视野相交、过滤小于
16 像素的八连通域。低阈值有助于保留与可信血管相连的弱响应，
也可能让假阳性沿弱背景传播。高阈值和连通域面积共同限制这一风险。

掩膜按行优先展开为 512 平方个二值元素，RLE 记录从 1 开始的游程起点与
正整数长度。生成 \code{Id,Predicted} 前，程序逐图解码并检查掩膜完全
相等。该步骤防止行列顺序、游程边界或空掩膜表示造成提交错误。

\section{本地训练与对照设计}
\label{sec:native-train}
\subsection{原生分辨率 U-Net}
原生路径在约 $584\times565$ 的图像上取绿色通道，做局部对比度增强
CLAHE，并计算 Frangi 管状响应。两通道分别提供亮度结构和血管形状线索。
训练在 FOV 内采样 $64\times64$ patch，并对图像与标签同步翻转。
每轮采样 8\,000 个 patch，16 张图用于训练，4 张图用于验证。
验证使用整幅补边到 608 的图像，避免以随机 patch 的分布代替整图评估。

U-Net 由卷积编码器、跳跃连接和上采样解码器构成，接单通道 logits 头。
两个损失项分别关注像素判别与前景重叠。按原生训练器配置，
\[
\mathcal L_{\rm native}
 =0.5\,\mathrm{BCEWithLogits}(\mathbf z,\mathbf g)
 +0.5\,\mathcal L_{\rm softDice}(\sigma(\mathbf z),\mathbf g).
\]
训练用 Adam 和余弦退火。日志保存了 40 个完整 epoch，验证 F1
在第 15 轮为 0.80780。历史总结将其称为 FOV Dice。现存训练器的验证
调用在补边后整图上累加像素计数，未向指标函数传入显式 FOV 掩膜，
因此本文将 0.80780 标为原生验证 Dice，并保留这项口径差异。
该日志属于 \code{e031_reti_native}。另一个
\code{e030_small_unet_frangi_f0} 运行的内部配置、日志和最佳轮次
分别归档，不能把其轮次或超参数并入原生模型。

\begin{table}[H]
\centering\small
\caption{分割主线和自训对照的配置证据}\label{tab:drive-cfg}
\begin{tabularx}{\linewidth}{@{}P{35mm}YY@{}}
\toprule
项目 & 冻结 FR-UNet & 原生 U-Net \\
\midrule
权重来源 & 上游 DRIVE epoch 40 & 课程训练图自训，验证选模 \\
输入 & $1\times512\times512$ 灰度 & CLAHE 与 Frangi 双通道 \\
训练采样 & 上游训练日志待确认 & $64\times64$，8\,000 patches/epoch \\
训练记录 & 课程路径加载权重推理 & 40 轮，最佳为第 15 轮 \\
损失与优化 & 上游设置按原论文核查 & BCE--Dice 各 0.5，Adam，余弦退火 \\
提交推理 & DTI、FOV 和面积过滤 & 原生概率图还原与 512 缩放 \\
\bottomrule
\end{tabularx}
\end{table}
\chart{drive_training.pdf}{原生 U-Net 的 40 轮训练与验证日志}{fig:drive-curve}
{验证值按训练器累加整幅补边图像的像素计数得到。标记点为第 15 轮的 0.80780。}

\subsection{对照目标}
保持同一 FR 权重，比较不同双阈值、单阈值或推理尺度，观察细血管召回
与背景误检之间的折衷。再将原生 U-Net 的概率和 FR 的概率按 25/75
加权，检验两条输入管线是否互补。自训模型只用课程训练标签，因此可以和
外部权重形成来源清晰的对照。不同架构、输入和评价范围使跨组差值难以
归结为单个因素。

\section{实验结果与错误分析}
表~\ref{tab:drive-results} 保留冻结配方与代表性失败实验。公开榜最佳
0.81766、私有榜 0.82534，均来自编号 \code{57946948}。
本地对全部 20 张课程训练图的 512 全图评估为 0.78712。上游权重的
训练域包含 DRIVE 训练图，该值用于核对实现输出。项目没有获得高于
0.855 的已验证公开榜结果，离该目标差 0.03734。

\begin{table}[H]
\centering\small
\caption{血管分割的 Kaggle 提交对照}\label{tab:drive-results}
\begin{tabularx}{\linewidth}{@{}Yrr@{}}
\toprule
配置 & Public Dice & Private Dice \\
\midrule
FR，DTI 0.65/0.40 & \textbf{0.81766} & \textbf{0.82534} \\
FR，DTI 0.80/0.40 & 0.81618 & 0.82399 \\
FR 0.75 与原生 U-Net 0.25 & 0.80897 & 0.81958 \\
FR，原生补边，八视图 TTA & 0.79524 & 0.80705 \\
原生 U-Net，阈值 0.35 & 0.78848 & 0.80285 \\
FR，单阈值 0.55，关闭 DTI & 0.77485 & 0.78705 \\
原生 U-Net，阈值 0.50 & 0.75079 & 待核查 \\
\bottomrule
\end{tabularx}
\end{table}
\chart{drive_results.pdf}{后处理与概率融合的分割结果}{fig:drive-results}
{左图使用实验日志中的公开榜结果。右图用平台提交记录对照公开榜与私有榜，两组坐标轴按所示范围显示。}

相对 0.65/0.40 配方，保留 DTI 且将高阈值提高到 0.80，公开榜
降低 0.00148。关闭 DTI 并用单阈值 0.55，下降 0.04281。
后者同时改变种子阈值和扩张规则，差值反映整套后处理变化。
原生模型在相同评估口径下由阈值 0.50 改为 0.35，公开榜提高 0.03769，
说明该模型对概率阈值敏感。仍有较多细小分支漏检的可能，
也可能受训练标签与提交尺寸之间的缩放影响。概率融合比原生模型更接近
冻结结果，较 FR 单模型仍低 0.00869。由此保留冻结提交及其哈希。

\ifLocalReport 图~\ref{fig:drive-case} 用冻结概率和 RLE 解码展示推理实例。眼底边缘
和细血管端点更容易出现低响应。测试样例无标注可用于判定其真实误差，
因此图中仅展示可见结构和模型输出。训练记录中的小样本验证也难以覆盖全部
成像条件。\fi
\privatechart{drive_case.png}{冻结 FR-UNet 的测试预测实例}{fig:drive-case}
{左侧为课程测试图，中间为存档血管概率，右侧为提交 RLE 解码的掩膜。三图对齐到提交尺寸。}

\section{本章小结}
最终配方把 FR-UNet 公开权重、双阈值扩张与提交格式检查串成可复查路径。
原生分辨率训练、不同阈值和概率融合提供了具体的负面结果。最好的公开与私有
Dice 分别为 0.81766 和 0.82534，结果适用于课程测试设置。
